\subsection{Reliability Analysis}

To assess the stability of the proposed evaluation framework, we measure inter-rater reliability using the intraclass correlation coefficient (ICC). ICC is computed separately at the level of problem-solution pairs $(p,s)$ treating multiple evaluation runs $x_r^{(p, s)}$ as different measurements. For each problem-solution pair $(p,s)$ we estimate a set of intraclass correlation coefficients:

\[
ICC = \{ICC(1,1), ICC(A,1), ICC(C,1), ICC(1,k), ICC(A,k), ICC(C,k)\}
\]

Thus, every problem-solution pair yields a different profile across multiple $ICC = \{ICC_{type}\}$ definitions:

\[
ICC^{(p,s)} = [ICC_{type}^{(p,s)}]
\]

To obtain dataset-level reliability estimates, we aggregate each ICC type independently across all pairs:

\[
\overline{ICC}_{type} = \frac{1}{|P|} \sum_{(p,s) \in P} ICC^{(p,s)}_{type}
\]

We additionally report standard deviations across pairs to capture heterogeneity in reliability across different tasks. 
The resulting aggregated reliability estimates across all problem-solution pairs are reported in Table \ref{table:reliability_analysis}.

\begin{table}[h]
\centering
\begin{tabular}{lcccccc}
\hline
 & ICC(1,1) & ICC(A,1) & ICC(C,1) & ICC(1,k) & ICC(A,k) & ICC(C,k) \\
\hline
Mean & 0.82 & 0.83 & 0.87 & 0.96 & 0.96 & 0.97 \\
Std  & 0.09 & 0.09 & 0.06 & 0.03 & 0.02 & 0.02 \\
\hline
\end{tabular}
\caption{Inter-rater reliability (ICC) of the proposed evaluation framework aggregated across all problem-solution pairs $(p,s)$. Values are reported as mean and standard deviation across pairs.}
\label{table:reliability_analysis}
\end{table}

The results indicate high inter-class reliability across all ICC formulations.
Single measure ICC values (ICC(1,1), ICC(A,1), ICC(C,1)) show significant agreement between individual evaluation runs, with values ranging from 0.82 to 0.87. 
Meanwhile average-measure ICC types (ICC(1,k), ICC(A,k), ICC(C,k)) are higher (0.96-0.97), indicating that aggregating multiple evaluation runs effectively reduces stochastic variance in LLM-based judgments and leads to highly stable aggregated scores. 
These results suggest that the proposed evaluation framework produces consistent assessments across repeated runs.

This confirms that variability in the evaluation is primarily driven by per-run stochasticity rather than systematic disagreement among evaluations. However, these estimates remain conditioned on the chosen LLM model and may vary under different judge models.